
\section{Applications to the Structure of Algebraic Groups}
\SGARefTarget{sga3:xiv:6}

\paragraph{\textbf{Theorem 6.1.}}
\SGARefTarget{sga3:xiv:6.1}
{\itshape Let \(G\) be a smooth algebraic group over a field
\(k\). Consider the scheme \(\mathcal T\) of maximal tori of
\(G\), isomorphic to the scheme \(\mathcal C\) of Cartan
subgroups of \(G\). This is a homogeneous space under \(G^0\) and
a smooth affine connected algebraic scheme
\textup{(\SGARefLink{sga3:xii:7.1}{XII~7.1}\textup{\SGARefLink{sga3:xii:7.1:item:d}{(d)}})}. Then \(\mathcal C\) is a rational
variety: its field of rational functions is a purely
transcendental extension of \(k\).}

We first prove the theorem when \(k\) is infinite. We may plainly
suppose that \(G\) is connected, since \(\mathcal T\), and hence
\(\mathcal C\), is unchanged when \(G\) is replaced by \(G^0\).
Moreover, by \SGARefLink{sga3:xii:7.6}{XII~7.6}, \(\mathcal C\) is unchanged when \(G\) is
divided by a central subgroup. We may therefore first divide by
the center of \(G\) and suppose that \(G\) is affine
\textup{(\SGARefLink{sga3:xii:6.1}{XII~6.1})}, and then divide by its reductive center
\textup{(\SGARefLink{sga3:xii:4.1}{XII~4.1} and \SGARefLink{sga3:xiv:4.4}{4.4})} and suppose that the reductive center
of \(G\) is trivial \textup{(\SGARefLink{sga3:xii:4.7}{XII~4.7}\textup{\SGARefLink{sga3:xii:4.7:item:b}{(b)}})}.

We proceed by induction on \(n=\dim G\), assuming the result for
dimensions \(n'<n\). If \(G\) is nilpotent, then \(\mathcal C\)
consists of a single \(k\)-rational point and the assertion is
trivial. Otherwise, the Lie algebra of \(G\) is not nilpotent
\textup{(\SGARefLink{sga3:xiv:1.3}{1.3})}. Its Cartan subalgebras, and therefore the
subgroups of type \((\mathrm C)\) of \(G\), have some dimension
\(n'<n\). Consider the morphism
\[
  \mathcal C\longrightarrow\mathcal D
\]
of \SGARefLink{sga3:xiv:5.5}{5.5}. By \SGARefLink{sga3:xiv:3.10}{3.10}, since the infinite-field Lie algebra
\(\mathfrak g\) contains a regular element, \(\mathcal D\) is a
rational variety. Thus its function field \(K\) is purely
transcendental over \(k\).

Let \(x\) be the generic point of \(\mathcal D\). By \SGARefLink{sga3:xiv:5.5}{5.5}, the
fiber of \(\mathcal C\) at \(x\) is the scheme of Cartan
subgroups of a smooth connected algebraic group \(D_x\) over
\(K=\kappa(x)\), namely the generic subgroup of type
\((\mathrm C)\) of \(G\). The function field \(L\) of
\(\mathcal C\) is therefore the function field of
\(\mathcal C_{D_x}\). Since \(\dim D_x=n'<n\), the induction
hypothesis makes \(L\) purely transcendental over \(K\), and
hence, by transitivity, purely transcendental over \(k\).

Suppose now that \(k\) is finite. A different argument is
required. We may again suppose that \(G\) is affine and
connected. Since \(k\) is perfect, the radical of
\(G_{\bar k}\) descends to a subgroup \(R\) of \(G\).
Suppose first that \(R\ne G\), so that \(G\) is not solvable, and
let
\[
  u\colon G\longrightarrow G'=G/R
\]
be the canonical morphism. By \SGARefLink{sga3:xii:7.1}{XII~7.1}\textup{\SGARefLink{sga3:xii:7.1:item:e}{(e)}}, it induces
\[
  v\colon\mathcal C_G\longrightarrow\mathcal C_{G'},
  \qquad C\longmapsto u(C).
\]
Let \(x\) be the generic point of \(\mathcal C_{G'}\), and put
\(K=\kappa(x)\). The fiber \(v^{-1}(x)\) is the scheme of Cartan
subgroups of \(G_K\) whose image in \(G'_K\) is the generic
Cartan subgroup \(C'_x\) of \(G'\). By \SGARefLink{sga3:xii:7.9}{XII~7.9}\textup{\SGARefLink{sga3:xii:7.9:item:c:02}{(c)}}, it is
also the scheme of Cartan subgroups of
\[
  H=u_K^{-1}(C'_x).
\]
The field \(K\) is infinite. The infinite-field case just proved
therefore shows that the function field \(L\) of
\(\mathcal C_G\), which is also that of \(\mathcal C_H\), is
purely transcendental over \(K\).\footnote{The French source
refers here to the ``already proved part of \SGARefLink{sga3:xiv:6.1}{7.1}''; the argument
is the infinite-field part of \SGARefLink{sga3:xiv:6.1}{Theorem 6.1}.}

It remains to prove that \(K\) is purely transcendental over
\(k\), and we are thus reduced to the case in which \(G\) is
semisimple. Dividing by its reductive center, we may further
suppose that \(G\) is adjoint. Then \SGARefLink{sga3:xiv:3.18}{3.18} gives
\(\mathcal C_G\simeq\mathcal D_G\), and by \SGARefLink{sga3:xiv:3.10}{3.10} it is enough to
show that \(\mathfrak g\) has a regular point. As noted in \SGARefLink{sga3:xiv:3.20}{3.20},
this is an unpublished result of Chevalley.

It remains only to treat the case in which \(k\) is finite and
\(G\) is connected, affine, and solvable. In fact, the next
result holds over an arbitrary field.
